% TOP_PROBLEM: {"id": 3347, "title": "Quartic rationally derived polynomials", "category_id": 1, "category": {"id": 1, "name": "number_theory", "display_name": "Number Theory", "description": "Properties of integers, prime numbers, Diophantine equations.", "slug": "number-theory", "order_index": 1, "created_at": "2026-07-31T15:26:25.670Z"}, "status": "open", "problem_number": "OPG-791", "set_id": 12}
% FIRST_POSTED: 2026-10-02
% ATTEMPT_STATUS: unresolved
% UnsolvedMath ID 3347; source catalogue code OPG-791
% !TeX program = lualatex
% Research attempt by GPT-6 Astra Ultra; no claim of a new complete solution.
\documentclass[11pt,a4paper]{article}
\usepackage[margin=25mm]{geometry}
\usepackage{fontspec}
\setmainfont{lmroman10-regular.otf}[BoldFont=lmroman10-bold.otf,
 ItalicFont=lmroman10-italic.otf,BoldItalicFont=lmroman10-bolditalic.otf]
\usepackage{amsmath,amssymb,mathtools}
\usepackage{unicode-math}
\setmathfont{latinmodern-math.otf}
\AtBeginDocument{\let\setminus\smallsetminus}
\usepackage{xurl,hyperref}
\hypersetup{colorlinks=true,urlcolor=blue}
\setcounter{secnumdepth}{0}
\setlength{\parindent}{0pt}
\setlength{\parskip}{5pt}
\setlength{\emergencystretch}{3em}
\begin{document}
\section*{Quartic rationally derived polynomials}\label{3347}
{\small UnsolvedMath ID 3347 · OPG-791 · Number Theory · OpenGarden}
\subsection{Definitions and mathematical statement}
A polynomial $p\in\mathbb Q[x]$ of degree four is rationally derived if
each nonconstant derivative $p^{(j)}$, for $0\le j\le3$, splits as a
product of linear factors in $\mathbb Q[x]$. Roots are counted with
multiplicity. Its fourth derivative is a nonzero constant and imposes
no condition; the identically zero later derivatives are excluded.
The conjecture is
\[
 \nexists p\in\mathbb Q[x],\quad \deg p=4,
 \quad p\text{ has four distinct roots},
 \quad p,p',p'',p^{(3)}\text{ split over }\mathbb Q.
\]
Equivalently, decide whether there are rational numbers
$a_1<a_2<a_3<a_4$ such that the derivatives of
$\prod_{i=1}^4(x-a_i)$ all have rational roots. Multiplying the
polynomial by a nonzero rational number does not affect the question.
An affine change $x\mapsto ux+v$, with $u,v\in\mathbb Q$ and $u\ne0$,
also preserves rational derivedness and distinctness. Consequently the
search can be normalized to roots $0,1,s,t$ with $1<s<t$ rational.
The rationality of all four original roots alone is insufficient:
the cubic and quadratic derivatives must split as well. A counterexample
must meet every one of these conditions simultaneously.
\subsection{Short English statement}
Can a quartic with four distinct rational roots also have only rational roots in every nonconstant derivative?
\subsection{Research attempt}
\paragraph{Scope and process.}
This is a GPT-6 Astra Ultra research attempt dated 2 October 2026, using
primary-source browsing and elementary mathematical checks. The canonical
definition above is retained. Results below are restricted results or
reductions, not a claim of a new solution to the entire catalogue question.
No formal proof assistant or external mathematical referee checked this text.


\paragraph{Source check and normalization.}
Flynn [R1] explicitly separates the conjecture about four distinct
quartic roots from the conjecture about a quintic having a triple root;
his theorem proves the latter. It must not be reported as a proof of
the former. The primary text also justifies rational affine changes
of variable. We use those changes to isolate a necessary conic and
exclude every symmetric root configuration. The literature search for
this attempt did not locate a proof settling the quartic target;
that search is not a substitute for one.

Scale the leading coefficient to one and translate by the rational
mean of the four roots. The result has the form
\[
                 p(x)=x^4-Ax^2+Bx+C,\qquad A,B,C\in\mathbb Q.
\]
If its four roots are distinct and real, Rolle's theorem gives three
distinct real roots of $p'$ and two of $p''$. Since
$p''(x)=12x^2-2A$, necessarily $A>0$. Requiring its roots to be rational
means $A=6t^2$ for some nonzero $t\in\mathbb Q$.

\paragraph{The remaining derivative condition.}
Write the rational roots of $p'/4$ as $r,s,-r-s$, using its zero
quadratic coefficient. Expanding the product shows that splitting of
the first two derivatives requires
\[
 r^2+rs+s^2=3t^2,\qquad B=4rs(r+s),\qquad A=6t^2.
\]
For a candidate with four distinct roots the three displayed derivative
roots must also be distinct. Conversely, any rational triple satisfying
these equations gives a polynomial whose positive-degree derivatives
split, whatever $C$ is. This converse deliberately says nothing about
the splitting of $p$ itself. The fourth and later derivatives add no
nonconstant condition, and the third derivative is $24x$.

\paragraph{Excluding symmetric roots.}
Suppose the original four rational roots are symmetric about some
centre. Their mean is that centre, so it is rational; after the
translation they are $u,-u,v,-v$, with $uv\ne0$ and $u^2\ne v^2$.
Then
\[
 p(x)=(x^2-u^2)(x^2-v^2),\qquad A=u^2+v^2>0,
\]
and $B=0$. The nonzero roots of $p'$ have square $A/2$, whereas the
roots of $p''$ have square $A/6$. If both derivatives split over
$\mathbb Q$, choosing a nonzero root from each would give two rational
numbers whose ratio has square three. That is impossible: writing a
putative square root of three as a reduced fraction forces both its
numerator and denominator divisible by three. Thus no quartic with
four distinct roots symmetric about a centre is rationally derived.
In particular this excludes every four-term rational arithmetic
progression. The nonzero condition on the selected roots is ensured
by $A>0$.

\paragraph{A checked obstruction to stopping at the conic.}
The derivative equations have many rational points that do not solve
the original question. Consider the explicit polynomial
\[
               p(x)=x^4-6x^2-\frac{1144}{343}x.
\]
Its first derivative satisfies
\[
 \frac{p'(x)}4=(x+11/7)(x+2/7)(x-13/7),
 \qquad p''(x)=12(x-1)(x+1),\qquad p^{(3)}(x)=24x.
\]
Nevertheless, putting $z=7x$ gives
\[
             2401p(z/7)=z(z^3-294z-1144).
\]
The cubic reduces modulo five to $z^3+z+1$. Its values at
$0,1,2,3,4$ are respectively $1,3,1,1,4$ modulo five, so it has no
root there. A rational root of a monic integer polynomial is an
integer and would reduce to such a root; consequently this cubic has
no rational root. The quartic therefore does not split over
$\mathbb Q$, even though every nonconstant positive-order derivative
does. This failure is not merely the absence of real roots: the
cubic has sign changes on $(-20,-10)$, $(-10,0)$ and $(0,20)$.
It consequently has three distinct real roots, none zero, so the
quartic has four distinct real roots. Certificate [R2] checks the
factorization, residues and signs with exact rational arithmetic.

\paragraph{Final status and first remaining gap.}
\textbf{UNRESOLVED.} All symmetric configurations are excluded, and
the derivative conditions are reduced to explicit rational equations.
The missing step is to impose four rational roots on the original
quartic simultaneously with those equations in the asymmetric case.
Solving the conic alone, or finding real roots numerically, does not
supply that step; the displayed example verifies the precise failure.

\subsection{Sources}
\begin{itemize}
\item[S1] ULAM.AI and the UnsolvedMath Contributors, supplied problems.json, record 3347 (OPG-791), OpenGarden collection. Catalogue identity and source transcription; CC BY 4.0.
\url{https://huggingface.co/datasets/ulamai/UnsolvedMath}.
\item[S2] Open Problem Garden, Quartic rationally derived polynomials. Original problem and discussion; the mathematical formulation below is expanded from the supplied source record.
\url{http://www.openproblemgarden.org/op/quartic_rationally_derived_polynomials}.

\item[R1] E. V. Flynn, \emph{On Q-derived polynomials}, Proceedings of
the Edinburgh Mathematical Society 44 (2001), 103--110. Author's
primary PDF, especially Conjectures 1.1 and 1.2, checked 2 October 2026.
\url{https://people.maths.ox.ac.uk/flynn/arts/art16.pdf}.
\item[R2] Reproducible exact checks:
\path{scripts/check_attempt_batch03_colouring_2026_10_02.py}, section 3347.

\end{itemize}
\end{document}
